% ---
% id: reductions/lem:top
% kind: lem
% title: "Flag PVM as a top element"
% summary: >-
%   Every POVM with a finite alphabet squashes to the flag PVM on its outcomes
%   through a measure-and-prepare channel, and to any PVM on the same alphabet
%   that vanishes only where the POVM does. A squasher therefore always
%   exists; the work lies in finding a target tight enough to leave a key rate.
% prominence: standard
% status: accepted
% uses: [def:squashing]
% assumes: []
% refs: []
% provenance: newly-derived
% lean:
% ---

A receiver is rarely characterized as a single POVM on a single space. It is
characterized \emph{conditionally}: below a photon-number cutoff we know what
it does and above it we do not; on one value of a classical branch register it
behaves as specified and on another we promise nothing. The lemma shows first
that existence of a squasher is never the difficulty. The work of this
module's constructions is therefore not to find a squasher but to find one
whose target is tight enough to leave a key rate.

\begin{lemma}\label{lem:top}
Let $M$ be any POVM on any $\Bsys$ with finite alphabet $\Vset$, and let
$\Ppvm=\{\dyad{v}\}_{v\in\Vset}$ be the \emph{flag PVM} on
$\Hs{\Flag}=\mathbb{C}^{|\Vset|}$ with orthonormal basis
$\{\ket{v}\}_{v\in\Vset}$. Then
\begin{equation}\label{eq:mpsquasher}
  M\;\dsq_{\Lmp}\;\Ppvm,
  \qquad
  \Lmp(\rho)\coloneqq\sum_{v\in\Vset}\tr[M_v\rho]\,\dyad{v} .
\end{equation}
More generally, if $P$ is any PVM on a space $\Qsys$ with the same alphabet and
$P_v=0\Rightarrow M_v=0$, then $M\dsq P$.
\end{lemma}

\begin{proof}
$\Lmp$ is a measure-and-prepare channel: CP, and trace preserving because
$\sum_v\tr[M_v\rho]=\tr\rho$. Its adjoint is
$\Lmp^{\dagger}(X)=\sum_v\bra{v}X\ket{v}\,M_v$, since
$\tr[X\Lmp(\rho)]=\sum_v\tr[M_v\rho]\bra{v}X\ket{v}$; evaluating at
$X=\dyad{v'}$ gives $M_{v'}$. For the general case pick, for each $v$ with
$M_v\neq0$, a state $\omega_v$ with $P_v\omega_vP_v=\omega_v$ (possible since
$P_v\neq0$), and set
$\Lambda(\rho)\coloneqq\sum_{v:M_v\neq0}\tr[M_v\rho]\,\omega_v$. This is trace
preserving because the omitted terms vanish and $\sum_vM_v=\idon{\Bsys}$, and
$\Lambda^{\dagger}(P_{v'})=\sum_{v:M_v\neq0}\tr[P_{v'}\omega_v]\,M_v=M_{v'}$,
because $\tr[P_{v'}\omega_v]=\delta_{vv'}$ by orthogonality of the ranges of a
PVM, and because the identity reads $0=0$ when $M_{v'}=0$.
\end{proof}
